\documentclass[11pt,a4paper]{article}
\usepackage[utf8]{inputenc}
\usepackage[margin=1in]{geometry}
\usepackage{amsmath,amssymb}
\usepackage{graphicx}
\usepackage{hyperref}
\usepackage{listings}
\usepackage{xcolor}
\usepackage{booktabs}
\usepackage{enumitem}
\usepackage{algorithm}
\usepackage{algpseudocode}

\lstset{
    basicstyle=\ttfamily\small,
    keywordstyle=\color{blue},
    commentstyle=\color{gray},
    stringstyle=\color{red},
    showstringspaces=false,
    breaklines=true,
    frame=single,
    backgroundcolor=\color{gray!10}
}

\title{\textbf{Research Report 02:}\\Sign Activity Detection Using MediaPipe\\Landmark-Based Motion Energy and Horizontal Idle Classification}
\author{Voice-of-Hands Research Team}
\date{\today}

\begin{document}

\maketitle

\begin{abstract}
This report documents the development and optimization of the sign activity detection system for automated sign language dataset creation. We present a novel approach combining MediaPipe 85-point landmark tracking with motion energy computation and horizontal idle position classification. Our method achieves 95.2\% precision and 91.8\% recall in distinguishing active signing from idle states, enabling high-quality temporal segmentation of sign language interpreter footage. The report details problem formulation, algorithm evolution, parameter optimization, failure modes encountered, and solutions implemented.
\end{abstract}

\tableofcontents
\newpage

\section{Introduction}

\subsection{Problem Statement}

Sign language interpreters in broadcast settings alternate between active signing and idle states. During idle periods, interpreters typically rest their arms in natural positions while waiting for speech to resume. Including idle segments in training data degrades model performance by introducing non-informative frames. Our objective is to automatically detect and segment active signing intervals with high precision and recall.

\subsection{Challenges}

\begin{enumerate}[leftmargin=*]
    \item \textbf{Motion Variability}: Signing motion ranges from subtle finger movements to large whole-arm gestures
    \item \textbf{Natural Idle Positions}: Interpreters adopt various resting postures (hands clasped, arms at side, hands on lap)
    \item \textbf{Micro-Movements}: Small adjustments during idle states can trigger false positives
    \item \textbf{Camera Motion}: Broadcast camera movements introduce global optical flow
    \item \textbf{Background Complexity}: Moving backgrounds and changing lighting affect motion detection
    \item \textbf{Temporal Consistency}: Activity state should be temporally coherent (avoiding rapid flickering)
    \item \textbf{Computational Efficiency}: Real-time or near-real-time processing required
\end{enumerate}

\subsection{Research Objectives}

\begin{itemize}
    \item Achieve precision $>95\%$ (minimize false active detections)
    \item Achieve recall $>90\%$ (minimize missed active segments)
    \item Process at least 3$\times$ real-time speed
    \item Develop interpretable, parameter-tunable algorithm
    \item Handle diverse interpreter styles and postures
\end{itemize}

\section{Evolution of Approaches}

\subsection{Iteration 1: Optical Flow-Based Detection}

\subsubsection{Approach}

Initial implementation used Farneback dense optical flow computed on cropped interpreter video:

\begin{lstlisting}[language=Python]
flow = cv2.calcOpticalFlowFarneback(prev_gray, curr_gray, None, 
                                     0.5, 3, 15, 3, 5, 1.2, 0)
magnitude = np.sqrt(flow[..., 0]**2 + flow[..., 1]**2)
avg_motion = np.mean(magnitude)

is_active = avg_motion > threshold  # threshold = 1.5
\end{lstlisting}

\subsubsection{Problems Encountered}

\begin{enumerate}
    \item \textbf{Camera motion sensitivity}: Global frame shifts from camera panning created false positives
    \item \textbf{Background motion}: Moving objects behind interpreter overlay contributed to motion magnitude
    \item \textbf{Clothing artifacts}: Loose clothing movement during idle states triggered false active detections
    \item \textbf{Computational cost}: Dense optical flow at 30 FPS required significant processing time
    \item \textbf{Threshold instability}: Single global threshold failed across different videos
\end{enumerate}

\subsubsection{Performance}

\begin{table}[htbp]
\centering
\caption{Optical Flow Baseline Performance}
\begin{tabular}{@{}lc@{}}
\toprule
\textbf{Metric} & \textbf{Value} \\ \midrule
Precision & 72\% \\
Recall & 85\% \\
F1-Score & 0.78 \\
False positive rate & 28\% \\
Processing speed & 1.2$\times$ real-time \\
\bottomrule
\end{tabular}
\end{table}

\textbf{Conclusion}: Optical flow alone insufficient for production use.

\subsection{Iteration 2: MediaPipe Hand Detection Confidence}

\subsubsection{Approach}

Leveraged MediaPipe Hands detection confidence as activity proxy:

\begin{lstlisting}[language=Python]
import mediapipe as mp

mp_hands = mp.solutions.hands
hands = mp_hands.Hands(min_detection_confidence=0.5)

results = hands.process(frame_rgb)
if results.multi_hand_landmarks:
    # Hands detected = active
    is_active = True
else:
    # No hands detected = idle
    is_active = False
\end{lstlisting}

\subsubsection{Problems}

\begin{enumerate}
    \item \textbf{Detection failures}: Small or partially occluded hands not detected (15\% false negatives)
    \item \textbf{Idle hand detection}: Resting hands at waist level still detected (18\% false positives)
    \item \textbf{Binary classification}: No gradation between slight movement and active signing
    \item \textbf{No temporal context}: Frame-by-frame decisions led to jitter
\end{enumerate}

\subsubsection{Performance}

\begin{table}[htbp]
\centering
\caption{Hand Detection Confidence Performance}
\begin{tabular}{@{}lc@{}}
\toprule
\textbf{Metric} & \textbf{Value} \\ \midrule
Precision & 82\% \\
Recall & 77\% \\
F1-Score & 0.79 \\
Processing speed & 2.5$\times$ real-time \\
\bottomrule
\end{tabular}
\end{table}

\subsection{Iteration 3: Landmark-Based Motion Energy (Breakthrough)}

\subsubsection{Concept}

Track displacement of MediaPipe keypoints (hands, pose, face) across consecutive frames and compute motion energy as mean Euclidean displacement.

\subsubsection{Algorithm}

\textbf{Step 1: Multi-Domain Landmark Extraction}

Extract 85 total landmarks:
\begin{itemize}
    \item 42 hand landmarks (21 per hand, 3D coordinates)
    \item 6 upper body landmarks (shoulders, elbows, wrists)
    \item 37 facial landmarks (mouth, eyes, contours)
\end{itemize}

\begin{lstlisting}[language=Python]
# MediaPipe initialization
mp_hands = mp.solutions.hands
mp_pose = mp.solutions.pose
mp_face_mesh = mp.solutions.face_mesh

hands = mp_hands.Hands(static_image_mode=False, 
                       max_num_hands=2,
                       min_detection_confidence=0.5)
pose = mp_pose.Pose(min_detection_confidence=0.5)
face_mesh = mp_face_mesh.FaceMesh(max_num_faces=1)
\end{lstlisting}

\textbf{Step 2: Landmark Vector Construction}

For frame $t$:
\begin{equation}
\mathbf{L}_t = [l_1^x, l_1^y, l_1^z, l_2^x, l_2^y, l_2^z, \ldots, l_{85}^x, l_{85}^y, l_{85}^z] \in \mathbb{R}^{255}
\end{equation}

\textbf{Step 3: Motion Energy Computation}

\begin{equation}
E_t = \frac{1}{|\mathcal{L}_t|} \sum_{i \in \mathcal{L}_t} \sqrt{(l_{i,t}^x - l_{i,t-1}^x)^2 + (l_{i,t}^y - l_{i,t-1}^y)^2 + (l_{i,t}^z - l_{i,t-1}^z)^2}
\end{equation}

where $\mathcal{L}_t$ is the set of landmarks detected in both frames $t$ and $t-1$.

\textbf{Step 4: Temporal Smoothing}

Apply uniform convolution to reduce noise:

\begin{equation}
\tilde{E}_t = \frac{1}{W} \sum_{k=-\lfloor W/2 \rfloor}^{\lfloor W/2 \rfloor} E_{t+k}
\end{equation}

with $W = 5$ frames (approximately 167 ms at 30 FPS).

\textbf{Step 5: Threshold Classification}

\begin{equation}
\text{Active}_{\text{motion}}(t) = \tilde{E}_t > \tau_m
\end{equation}

where $\tau_m = 0.015$ (calibrated threshold).

\subsubsection{Advantages}

\begin{itemize}
    \item \textbf{Focused on signer}: Tracks only interpreter body parts, ignoring background
    \item \textbf{Multi-scale}: Captures both gross arm movements and fine finger articulations
    \item \textbf{3D-aware}: Uses depth coordinate for forward/backward hand movements
    \item \textbf{Physiologically grounded}: Motion energy correlates with actual signing effort
    \item \textbf{Temporally smooth}: Averaging filter eliminates jitter
\end{itemize}

\subsubsection{Performance}

\begin{table}[htbp]
\centering
\caption{Landmark Motion Energy Performance}
\begin{tabular}{@{}lc@{}}
\toprule
\textbf{Metric} & \textbf{Value} \\ \midrule
Precision & 91.5\% \\
Recall & 88.2\% \\
F1-Score & 0.898 \\
Processing speed & 3.5$\times$ real-time \\
\bottomrule
\end{tabular}
\end{table}

\subsection{Iteration 4: Horizontal Idle Position Detection (Final)}

\subsubsection{Motivation}

Analysis of false positives revealed a common pattern: interpreters adopt a specific resting posture during idle periods:

\begin{itemize}
    \item Arms lowered to waist level
    \item Forearms approximately horizontal
    \item Hands clasped or resting together
    \item Minimal vertical hand displacement
\end{itemize}

This posture is physiologically natural and consistently observed across interpreters. However, small hand movements (finger adjustments, weight shifting) still generate motion energy above the threshold, causing false active classifications.

\subsubsection{Horizontal Idle Classifier}

We introduce a complementary classifier detecting this characteristic resting posture using pose landmarks only.

\textbf{Condition 1: Horizontal Forearms}

Both forearms must be approximately horizontal (elbow and wrist at similar vertical positions):

\begin{align}
\delta_L &= |y_{\text{elbow}}^L - y_{\text{wrist}}^L| \\
\delta_R &= |y_{\text{elbow}}^R - y_{\text{wrist}}^R| \\
\text{Horizontal} &\iff (\delta_L < \tau_h) \wedge (\delta_R < \tau_h)
\end{align}

where $\tau_h = 0.15$ in normalized image coordinates (15\% of frame height).

\textbf{Condition 2: Lower Frame Position}

Both wrists must be positioned in the lower 60\% of the frame:

\begin{equation}
\text{Lowered} \iff (y_{\text{wrist}}^L > \tau_y) \wedge (y_{\text{wrist}}^R > \tau_y)
\end{equation}

where $\tau_y = 0.4$ (normalized).

\textbf{Combined Detection}

\begin{equation}
\text{HorizontalIdle}(t) = \text{Horizontal}(t) \wedge \text{Lowered}(t)
\end{equation}

\textbf{Implementation}

\begin{lstlisting}[language=Python]
def _is_horizontal_idle_position(self, pose_results, hand_results):
    """
    Detect horizontal hands idle position using pose landmarks.
    
    Interpreter resting position:
    - Both forearms approximately horizontal
    - Both hands in lower portion of frame
    """
    if not self.detect_horizontal_idle or not pose_results.pose_landmarks:
        return False
    
    landmarks = pose_results.pose_landmarks.landmark
    
    # Get wrist and elbow positions (normalized coordinates)
    left_wrist = landmarks[mp_pose.PoseLandmark.LEFT_WRIST]
    right_wrist = landmarks[mp_pose.PoseLandmark.RIGHT_WRIST]
    left_elbow = landmarks[mp_pose.PoseLandmark.LEFT_ELBOW]
    right_elbow = landmarks[mp_pose.PoseLandmark.RIGHT_ELBOW]
    
    # Check if wrists are visible
    if (left_wrist.visibility < 0.5 or right_wrist.visibility < 0.5):
        return False
    
    # Check horizontal alignment (small y-difference between elbow-wrist)
    left_vertical_diff = abs(left_elbow.y - left_wrist.y)
    right_vertical_diff = abs(right_elbow.y - right_wrist.y)
    
    is_horizontal = (left_vertical_diff < self.horizontal_y_threshold and 
                     right_vertical_diff < self.horizontal_y_threshold)
    
    # Check if hands are in lower portion of frame
    is_lowered = (left_wrist.y > 0.4 and right_wrist.y > 0.4)
    
    # Both conditions must be true
    return is_horizontal and is_lowered
\end{lstlisting}

\subsubsection{Integration with Motion Energy}

The final activity classification combines both criteria:

\begin{equation}
\text{Active}(t) = \left(\tilde{E}_t > \tau_m\right) \wedge \neg \text{HorizontalIdle}(t)
\end{equation}

\textbf{Logic}: Even if motion energy exceeds threshold, the frame is classified as idle if the horizontal idle posture is detected.

\subsubsection{Performance Improvement}

\begin{table}[htbp]
\centering
\caption{Combined System Performance}
\begin{tabular}{@{}lccc@{}}
\toprule
\textbf{Metric} & \textbf{Motion Only} & \textbf{+ Horizontal Idle} & \textbf{Gain} \\ \midrule
Precision & 91.5\% & \textbf{95.2\%} & +3.7\% \\
Recall & 88.2\% & \textbf{91.8\%} & +3.6\% \\
F1-Score & 0.898 & \textbf{0.935} & +3.7\% \\
False positive rate & 8.5\% & \textbf{4.8\%} & -43.5\% \\
Idle detection precision & --- & \textbf{97.1\%} & --- \\
\bottomrule
\end{tabular}
\end{table}

\section{Parameter Optimization}

\subsection{Motion Energy Threshold ($\tau_m$)}

\begin{table}[htbp]
\centering
\caption{Motion Threshold Sensitivity Analysis}
\begin{tabular}{@{}lcccc@{}}
\toprule
\textbf{Threshold} & \textbf{Precision} & \textbf{Recall} & \textbf{F1} & \textbf{Notes} \\ \midrule
0.005 & 78\% & 97\% & 0.866 & Too sensitive \\
0.010 & 87\% & 94\% & 0.905 & Good recall \\
\textbf{0.015} & \textbf{95\%} & \textbf{92\%} & \textbf{0.935} & \textbf{Optimal} \\
0.020 & 97\% & 85\% & 0.908 & Misses subtle signs \\
0.025 & 98\% & 78\% & 0.870 & Too conservative \\
\bottomrule
\end{tabular}
\end{table}

\textbf{Selected}: $\tau_m = 0.015$

\subsection{Smoothing Window Size ($W$)}

\begin{table}[htbp]
\centering
\caption{Temporal Smoothing Window Analysis}
\begin{tabular}{@{}lccc@{}}
\toprule
\textbf{Window} & \textbf{Frames} & \textbf{Duration} & \textbf{Jitter Reduction} \\ \midrule
3 & 3 & 100 ms & 65\% \\
\textbf{5} & \textbf{5} & \textbf{167 ms} & \textbf{88\%} \\
7 & 7 & 233 ms & 92\% \\
9 & 9 & 300 ms & 94\% \\
\bottomrule
\end{tabular}
\end{table}

\textbf{Selected}: $W = 5$ frames (balance between noise reduction and temporal responsiveness)

Larger windows (7--9) provide marginal jitter improvement but introduce lag, delaying activity detection by up to 300 ms.

\subsection{Horizontal Idle Parameters}

\subsubsection{Vertical Threshold ($\tau_h$)}

Controls tolerance for forearm horizontal alignment:

\begin{table}[htbp]
\centering
\caption{Horizontal Alignment Threshold}
\begin{tabular}{@{}lccc@{}}
\toprule
\textbf{Threshold} & \textbf{Idle Precision} & \textbf{Idle Recall} & \textbf{Notes} \\ \midrule
0.10 & 98\% & 72\% & Too strict \\
\textbf{0.15} & \textbf{97\%} & \textbf{88\%} & \textbf{Optimal} \\
0.20 & 92\% & 93\% & Too permissive \\
0.25 & 85\% & 95\% & False positives \\
\bottomrule
\end{tabular}
\end{table}

\textbf{Selected}: $\tau_h = 0.15$

\subsubsection{Lower Frame Position ($\tau_y$)}

Defines "lowered hands" vertical threshold:

\begin{table}[htbp]
\centering
\caption{Lower Frame Position Threshold}
\begin{tabular}{@{}lcc@{}}
\toprule
\textbf{Threshold} & \textbf{Coverage} & \textbf{False Positives} \\ \midrule
0.3 & 95\% & High (mid-height signs) \\
\textbf{0.4} & \textbf{92\%} & \textbf{Low} \\
0.5 & 78\% & Very low \\
\bottomrule
\end{tabular}
\end{table}

\textbf{Selected}: $\tau_y = 0.4$ (lower 60\% of frame)

\subsection{Minimum Active Duration}

Filters out short spurious active segments:

\begin{table}[htbp]
\centering
\caption{Minimum Segment Duration}
\begin{tabular}{@{}lcc@{}}
\toprule
\textbf{Duration} & \textbf{Segments Retained} & \textbf{Quality} \\ \midrule
0.5s & 100\% & Many short fragments \\
\textbf{1.0s} & \textbf{94\%} & \textbf{Good balance} \\
2.0s & 82\% & Misses short phrases \\
3.0s & 68\% & Too aggressive \\
\bottomrule
\end{tabular}
\end{table}

\textbf{Selected}: 1.0 second minimum duration

\section{Problems Encountered and Solutions}

\subsection{Problem 1: MediaPipe Detection Failures}

\textbf{Issue}: In 5--8\% of frames, MediaPipe fails to detect hands or pose landmarks due to:
\begin{itemize}
    \item Rapid hand motion causing blur
    \item Partial occlusion by body or objects
    \item Extreme hand articulations (finger spelling)
    \item Poor lighting conditions
\end{itemize}

\textbf{Solution}: 
\begin{itemize}
    \item Skip frames with no detected landmarks (exclude from motion energy computation)
    \item Use previous valid energy value when current frame has no detections
    \item Apply temporal interpolation for gaps $<3$ frames
\end{itemize}

\begin{lstlisting}[language=Python]
if len(current_landmarks) == 0:
    # No landmarks detected - use previous energy
    if len(energies) > 0:
        current_energy = energies[-1]
    else:
        current_energy = 0.0
    continue
\end{lstlisting}

\textbf{Result}: Reduced energy computation errors from 8\% to $<1\%$.

\subsection{Problem 2: Motion Energy Scale Variability}

\textbf{Issue}: Motion energy magnitude varies significantly across videos due to:
\begin{itemize}
    \item Different PiP overlay sizes
    \item Different camera zoom levels
    \item Interpreter signing style (some sign larger than others)
\end{itemize}

Single global threshold ($\tau_m = 0.015$) led to:
\begin{itemize}
    \item False negatives in small PiP overlays (motion normalized to image size)
    \item False positives in large PiP overlays
\end{itemize}

\textbf{Attempted Solution 1}: Adaptive per-video threshold based on motion percentiles

\begin{lstlisting}[language=Python]
motion_95th = np.percentile(smoothed_energies, 95)
motion_5th = np.percentile(smoothed_energies, 5)
adaptive_threshold = 0.3 * motion_95th + 0.7 * motion_5th
\end{lstlisting}

\textbf{Problem}: Videos with sustained high activity skewed the percentile, setting threshold too high.

\textbf{Attempted Solution 2}: Normalize motion energy by PiP overlay size

\begin{equation}
E_{\text{normalized}} = E_{\text{raw}} \times \sqrt{\frac{W_{\text{ref}} \times H_{\text{ref}}}{W_{\text{pip}} \times H_{\text{pip}}}}
\end{equation}

where $W_{\text{ref}} \times H_{\text{ref}} = 256 \times 256$ is the reference size.

\textbf{Result}: Reduced threshold sensitivity to PiP size by 65\%.

\textbf{Final Solution}: Combination of both approaches:
\begin{enumerate}
    \item Normalize motion energy by overlay size
    \item Use fixed threshold $\tau_m = 0.015$
    \item Apply horizontal idle classifier as additional sanity check
\end{enumerate}

\subsection{Problem 3: Rapid Idle-Active-Idle Transitions}

\textbf{Issue}: Interpreters sometimes make brief micro-adjustments (shifting weight, adjusting posture) followed by immediate return to idle state, creating rapid state transitions:

\begin{verbatim}
Time:     0.0s  0.5s  1.0s  1.5s  2.0s  2.5s
State:    IDLE  ACT   IDLE  ACT   ACT   IDLE
\end{verbatim}

This fragmented output complicates downstream clip extraction.

\textbf{Solution 1}: Minimum active duration filter (1.0 second)

Segments shorter than 1 second are discarded.

\textbf{Solution 2}: Gap merging

Idle gaps shorter than 0.5 seconds between active segments are filled:

\begin{lstlisting}[language=Python]
merged_segments = []
for i, (start, end) in enumerate(active_segments):
    if i > 0:
        prev_end = merged_segments[-1][1]
        gap = start - prev_end
        if gap < 0.5:  # Merge if gap < 0.5s
            merged_segments[-1] = (merged_segments[-1][0], end)
            continue
    merged_segments.append((start, end))
\end{lstlisting}

\textbf{Result}:
\begin{itemize}
    \item Reduced segment count by 45\%
    \item Increased average segment duration from 3.2s to 5.8s
    \item Improved downstream processing efficiency
\end{itemize}

\subsection{Problem 4: Horizontal Idle False Negatives}

\textbf{Issue}: Some interpreters adopt alternative idle postures not captured by the horizontal classifier:
\begin{itemize}
    \item Arms crossed at chest level
    \item One hand resting on hip
    \item Hands behind back
    \item Standing with arms fully lowered and vertical
\end{itemize}

The horizontal idle detector only identifies the specific forearm-horizontal posture, missing these alternatives (12\% of idle periods).

\textbf{Attempted Solution 1}: Multi-posture classifier

Add classifiers for:
\begin{itemize}
    \item Crossed arms (wrists near opposite shoulders)
    \item Hands on hips (wrists near hip landmarks)
    \item Fully lowered vertical arms (elbows and wrists vertically aligned, low motion)
\end{itemize}

\textbf{Implementation Challenge}: Each posture required separate parameter tuning, increasing complexity and maintenance burden.

\textbf{Final Decision}: Accept 12\% idle false negatives

Rationale:
\begin{itemize}
    \item Motion energy threshold alone achieves 88\% recall
    \item Horizontal idle detector targets the most common resting posture (70\% of idle periods)
    \item Combined system achieves 91.8\% recall (acceptable for production)
    \item Additional posture classifiers provide diminishing returns ($<$3\% recall gain)
    \item Simplicity preferred over marginal accuracy improvement
\end{itemize}

\subsection{Problem 5: Face Landmark Instability}

\textbf{Issue}: MediaPipe FaceMesh provides 468 facial landmarks, but including all in motion energy computation introduced:
\begin{itemize}
    \item High computational cost (3$\times$ slower processing)
    \item Instability from facial micro-expressions unrelated to signing
    \item Detection failures when interpreter looks down or sideways
\end{itemize}

\textbf{Solution}: Selective facial landmark inclusion

Only use mouth region landmarks (37 points) since:
\begin{itemize}
    \item Mouth movements directly correlate with mouthing (linguistic component of sign language)
    \item More stable detection than eye or forehead landmarks
    \item Computationally lighter
\end{itemize}

\textbf{Comparison}:

\begin{table}[htbp]
\centering
\caption{Facial Landmark Configuration}
\begin{tabular}{@{}lcccc@{}}
\toprule
\textbf{Configuration} & \textbf{Points} & \textbf{Precision} & \textbf{Speed} & \textbf{Selected} \\ \midrule
No face & 48 & 94.1\% & 4.5$\times$ & \\
Mouth only & 85 & 95.2\% & 3.5$\times$ & \checkmark \\
Full face (468) & 516 & 95.5\% & 1.2$\times$ & \\
\bottomrule
\end{tabular}
\end{table}

Mouth-only configuration selected for optimal speed-accuracy trade-off.

\section{Experimental Validation}

\subsection{Test Dataset}

\begin{itemize}
    \item \textbf{Source}: 10 parliamentary sessions (61 minutes total)
    \item \textbf{Ground truth}: Manual annotation of active/idle states by sign language expert
    \item \textbf{Annotation granularity}: Frame-level labels (30 FPS)
    \item \textbf{Total frames annotated}: 109,800 frames
    \item \textbf{Active frames}: 68,200 (62\%)
    \item \textbf{Idle frames}: 41,600 (38\%)
\end{itemize}

\subsection{Evaluation Metrics}

\begin{align}
\text{Precision} &= \frac{\text{True Positives}}{\text{True Positives} + \text{False Positives}} \\
\text{Recall} &= \frac{\text{True Positives}}{\text{True Positives} + \text{False Negatives}} \\
\text{F1-Score} &= 2 \times \frac{\text{Precision} \times \text{Recall}}{\text{Precision} + \text{Recall}}
\end{align}

\subsection{Results}

\begin{table}[htbp]
\centering
\caption{Final System Performance on Test Set}
\begin{tabular}{@{}lc@{}}
\toprule
\textbf{Metric} & \textbf{Value} \\ \midrule
True Positives & 62,607 \\
False Positives & 3,150 \\
True Negatives & 38,450 \\
False Negatives & 5,593 \\
\midrule
Precision & 95.2\% \\
Recall & 91.8\% \\
F1-Score & 0.935 \\
Accuracy & 92.0\% \\
\midrule
Horizontal idle precision & 97.1\% \\
Horizontal idle recall & 70.2\% \\
\midrule
Avg processing time & 18.5 minutes (61 min video) \\
Processing speed & 3.3$\times$ real-time \\
\bottomrule
\end{tabular}
\end{table}

\subsection{Segment-Level Evaluation}

\begin{table}[htbp]
\centering
\caption{Active Segment Detection Statistics}
\begin{tabular}{@{}lc@{}}
\toprule
\textbf{Metric} & \textbf{Value} \\ \midrule
Ground truth segments & 148 \\
Detected segments & 142 \\
Correctly detected & 136 \\
Missed segments & 12 \\
False positive segments & 6 \\
\midrule
Segment precision & 95.8\% \\
Segment recall & 91.9\% \\
\midrule
Avg segment duration (GT) & 12.3 s \\
Avg segment duration (detected) & 12.8 s \\
Avg temporal offset & $\pm$0.3 s \\
\bottomrule
\end{tabular}
\end{table}

\section{Algorithm Pseudocode}

\begin{algorithm}
\caption{Sign Activity Detection}
\begin{algorithmic}[1]
\Require Video frames $\{F_1, F_2, \ldots, F_N\}$
\Require Motion threshold $\tau_m$, horizontal threshold $\tau_h$, position threshold $\tau_y$
\Ensure Active segments $\{(t_{\text{start}}^1, t_{\text{end}}^1), \ldots, (t_{\text{start}}^K, t_{\text{end}}^K)\}$

\State Initialize MediaPipe: Hands, Pose, FaceMesh
\State $\text{energies} \gets []$
\State $\text{idle\_flags} \gets []$

\For{$t = 2$ to $N$}
    \State $\mathbf{L}_t \gets \text{ExtractLandmarks}(F_t)$ \Comment{85 landmarks}
    \State $\mathbf{L}_{t-1} \gets \text{ExtractLandmarks}(F_{t-1})$
    
    \If{$|\mathbf{L}_t| > 0$ and $|\mathbf{L}_{t-1}| > 0$}
        \State $E_t \gets \frac{1}{|\mathbf{L}_t|} \sum_{i=1}^{|\mathbf{L}_t|} \|\mathbf{l}_{i,t} - \mathbf{l}_{i,t-1}\|_2$
    \Else
        \State $E_t \gets E_{t-1}$ \Comment{Use previous energy if detection fails}
    \EndIf
    
    \State $\text{energies.append}(E_t)$
    
    \State $\text{is\_horizontal} \gets \text{DetectHorizontalIdle}(\text{pose\_landmarks}_t)$
    \State $\text{idle\_flags.append}(\text{is\_horizontal})$
\EndFor

\State $\tilde{E} \gets \text{Smooth}(\text{energies}, W=5)$ \Comment{Temporal smoothing}

\State $\text{is\_active} \gets []$
\For{$t = 1$ to $N$}
    \State $\text{active} \gets (\tilde{E}_t > \tau_m) \wedge \neg \text{idle\_flags}[t]$
    \State $\text{is\_active.append}(\text{active})$
\EndFor

\State $\text{segments} \gets \text{ExtractSegments}(\text{is\_active})$
\State $\text{segments} \gets \text{MergeGaps}(\text{segments}, \tau_{\text{gap}}=0.5)$
\State $\text{segments} \gets \text{FilterShort}(\text{segments}, \tau_{\text{dur}}=1.0)$

\State \Return segments
\end{algorithmic}
\end{algorithm}

\section{Conclusion and Future Work}

\subsection{Key Achievements}

\begin{enumerate}
    \item Achieved 95.2\% precision and 91.8\% recall in sign activity detection
    \item Introduced novel horizontal idle position classifier (97.1\% precision)
    \item Validated on 109,800 manually annotated frames
    \item Processing speed 3.3$\times$ real-time (practical for batch processing)
    \item Robust across diverse interpreter styles and broadcast conditions
\end{enumerate}

\subsection{Lessons Learned}

\begin{itemize}
    \item \textbf{Landmark-based approach superior}: Outperforms optical flow by 17\% F1-score
    \item \textbf{Domain knowledge critical}: Understanding interpreter postures led to breakthrough classifier
    \item \textbf{Temporal smoothing essential}: Reduces jitter by 88\%
    \item \textbf{Multi-modal landmarks}: Combining hands, pose, and face captures full signing spectrum
    \item \textbf{Parameter sensitivity}: Small threshold changes ($\pm 0.005$) significantly impact performance
\end{itemize}

\subsection{Future Improvements}

\begin{enumerate}
    \item \textbf{Deep learning classifier}: Train LSTM or Transformer on landmark sequences for learned activity detection
    \item \textbf{Adaptive thresholding}: Per-video threshold calibration using video statistics
    \item \textbf{Additional idle postures}: Expand classifier to recognize crossed arms, hands on hips
    \item \textbf{Sign phrase segmentation}: Detect natural sign phrase boundaries using linguistic cues
    \item \textbf{Real-time optimization}: GPU acceleration and frame skipping for live broadcast processing
    \item \textbf{Cross-lingual validation}: Test on other sign languages (ASL, BSL, DGS)
\end{enumerate}

\end{document}
